\documentclass[../main.tex]{subfiles}

\begin{document}

\section*{第1問}

\begin{proof}[解]
$a_n=\displaystyle\int_1^e(\log x)^n\,dx$（$n=1,2,\dots$）とする。

\textbf{(1)} $n\ge3$の時，部分積分を用いる。まず
\[
a_{n-1}=\int_1^e(\log x)^{n-1}dx=\bigl[x(\log x)^{n-1}\bigr]_1^e-\int_1^ex\cdot(n-1)(\log x)^{n-2}\cdot\frac1x\,dx=e-(n-1)a_{n-2}
\]
（$n-1\ge2$の場合の同じ部分積分公式）。同様に
\[
a_n=\bigl[x(\log x)^n\bigr]_1^e-\int_1^ex\cdot n(\log x)^{n-1}\cdot\frac1x\,dx=e-n\,a_{n-1}
\]
この2式から$e$を消去すると，$e=a_{n-1}+(n-1)a_{n-2}$を$a_n=e-na_{n-1}$に代入して
\[
a_n=\{a_{n-1}+(n-1)a_{n-2}\}-na_{n-1}=(n-1)a_{n-2}-(n-1)a_{n-1}=(n-1)(a_{n-2}-a_{n-1})
\]

\textbf{(2)} $[1,e]$では$0\le\log x\le1$（②）。したがって$0\le(\log x)^{n+1}\le(\log x)^n$（$0\le t\le1$なら$t^{n+1}\le t^n$）。同区間で積分して
\[
0\le a_{n+1}\le a_n
\]
さらに，②の等号は$x=1$（両辺$0$）または$x=e$（両辺$1$）でしか成立しないから，$(1,e)$上では真の不等式$(\log x)^{n+1}<(\log x)^n$が成り立ち，正の測度の区間で成り立つので積分の不等式も真に：
\[
0<a_{n+1}<a_n
\]

\textbf{(3)} $n\ge2$とする。(1)（$2n\ge4\ge3$に適用）から
\[
a_{2n}=(2n-1)(a_{2n-2}-a_{2n-1})
\]
(2)から$a_{2n}<a_{2n-1}$，すなわち$-a_{2n-1}<-a_{2n}$だから，$a_{2n-2}-a_{2n-1}<a_{2n-2}-a_{2n}$。$(2n-1)>0$をかけて
\[
a_{2n}<(2n-1)(a_{2n-2}-a_{2n})
\]
整理すると$a_{2n}+(2n-1)a_{2n}<(2n-1)a_{2n-2}$，すなわち
\[
2n\,a_{2n}<(2n-1)a_{2n-2} \quad\therefore\ a_{2n}<\frac{2n-1}{2n}a_{2n-2}
\]
これをくり返し用いて（$n,n-1,\dots,2$について）
\[
a_{2n}<\frac{2n-1}{2n}\cdot\frac{2n-3}{2(n-1)}\cdots\frac34\,a_2
\]
$a_2$を求める。$\int(\log x)^2dx=x(\log x)^2-2x\log x+2x+C$（微分すれば確認できる）だから
\[
a_2=\bigl[x(\log x)^2-2x\log x+2x\bigr]_1^e=(e-2e+2e)-(0-0+2)=e-2
\]
以上から
\[
a_{2n}<\frac{(2n-1)(2n-3)\cdots3}{(2n)(2n-2)\cdots4}(e-2)=\frac{3\cdot5\cdots(2n-1)}{4\cdot6\cdots(2n)}(e-2)
\]
\end{proof}

\newpage

\section*{第2問}

\begin{proof}[解]
$\alpha$（1回目），$\beta$（2回目）は独立に$1,\dots,6$の一様分布，$E(\alpha)=E(\beta)=7/2$，$E(\alpha^2)=E(\beta^2)=\dfrac16\displaystyle\sum_{k=1}^6k^2=\dfrac{91}6$。$f(x)=(x-\alpha)(x-\beta)=x^2+sx+t$だから，解と係数の関係より
\[
s=-(\alpha+\beta),\qquad t=\alpha\beta
\]
また$f(x)^2=(x^2+sx+t)^2=x^4+2sx^3+(s^2+2t)x^2+2stx+t^2$だから
\[
a=2s,\qquad b=s^2+2t,\qquad c=2st,\qquad d=t^2
\]

\textbf{(1)} $\alpha,\beta$は独立だから
\[
E(s)=-E(\alpha)-E(\beta)=-7,\qquad E(t)=E(\alpha)E(\beta)=\frac{49}4
\]

\textbf{(2)}
\[
E(a)=2E(s)=-14
\]

$b=s^2+2t=(\alpha+\beta)^2+2\alpha\beta=\alpha^2+4\alpha\beta+\beta^2$だから，独立性より
\[
E(b)=E(\alpha^2)+4E(\alpha)E(\beta)+E(\beta^2)=2\cdot\frac{91}6+4\cdot\frac{49}4=\frac{91}3+49=\frac{238}3
\]

$c=2st=-2(\alpha+\beta)\alpha\beta=-2(\alpha^2\beta+\alpha\beta^2)$だから
\[
E(c)=-2\bigl\{E(\alpha^2)E(\beta)+E(\alpha)E(\beta^2)\bigr\}=-2\cdot2\cdot\frac{91}6\cdot\frac72=-\frac{637}3
\]

$d=t^2=\alpha^2\beta^2$だから
\[
E(d)=E(\alpha^2)E(\beta^2)=\left(\frac{91}6\right)^2=\frac{8281}{36}
\]
\end{proof}

\newpage

\section*{第3問}

\begin{proof}[解]
$D$の中心は$C:x^2+y^2=1$上にあるから，$C=\cos\theta,\ S=\sin\theta$として$P(C,S,0)$（$0\le\theta<2\pi$）とかける。$D$のある平面は常に$(0,1,0)$と直交するので，平面$y=S$内にあり
\[
D:\ (x-C)^2+z^2\le1,\quad y=S
\]

平面$y=k$（$-1\le k\le1$）による断面を考える。$\sin\theta=k$をみたす$\theta\in[0,\pi]$に対応して$\cos\theta=\pm\sqrt{1-k^2}$の2値をとり得るが，対称性からどちらも同じ形の断面を与える。$c=\sqrt{1-k^2}\ (=|\cos\theta|)$とすると，断面は$xz$平面内の2つの単位円板（中心$(\pm c,0)$）の和集合であり，この面積を$V(\theta)$とおく（$c=\cos\theta$，$0\le\theta\le\pi/2$として良い）。

2円板の共通部分（レンズ形）の面積は，対称性から
\[
4\int_c^1\sqrt{1-x^2}\,dx
\]
に等しいから
\[
V(\theta)=2\pi-4\int_c^1\sqrt{1-x^2}\,dx
\]
$x=\cos\alpha$（$\alpha:\theta\to0$として$x:c\to1$）とおくと，$dx=-\sin\alpha\,d\alpha$，$\sqrt{1-x^2}=\sin\alpha$（$0\le\alpha\le\pi/2$）だから
\[
\int_c^1\sqrt{1-x^2}\,dx=\int_0^\theta\sin^2\alpha\,d\alpha=\left[\frac\alpha2-\frac{\sin2\alpha}4\right]_0^\theta=\frac\theta2-\frac{\sin2\theta}4
\]
したがって
\[
V(\theta)=2\pi-4\left(\frac\theta2-\frac{\sin2\theta}4\right)=2\pi+\sin2\theta-2\theta
\]

もとめる体積$V$は，$y=\sin\theta$の対称性（$0\le\theta\le\pi/2$で$y:0\to1$，残りは対称）から
\[
\frac V2=\int_0^1V(\theta)\,dy=\int_0^{\pi/2}V(\theta)\cdot\frac{dy}{d\theta}\,d\theta=\int_0^{\pi/2}(2\pi+\sin2\theta-2\theta)\cos\theta\,d\theta
\]
\[
=\int_0^{\pi/2}\bigl(2\pi C+C\sin2\theta-2C\theta\bigr)d\theta \quad(C=\cos\theta) \quad\cdots\text{①}
\]
各項を計算する。
\[
\int_0^{\pi/2}C\,d\theta=\bigl[\sin\theta\bigr]_0^{\pi/2}=1
\]
\[
\int_0^{\pi/2}C\theta\,d\theta=\bigl[\theta\sin\theta+\cos\theta\bigr]_0^{\pi/2}=\left(\frac\pi2+0\right)-(0+1)=\frac\pi2-1
\]
\[
\int_0^{\pi/2}C\sin2\theta\,d\theta=2\int_0^{\pi/2}\sin\theta\cos^2\theta\,d\theta=2\left[-\frac{\cos^3\theta}3\right]_0^{\pi/2}=2\cdot\frac13=\frac23
\]
①に代入して
\[
\frac V2=2\pi\cdot1+\frac23-2\left(\frac\pi2-1\right)=2\pi+\frac23-\pi+2=\pi+\frac83
\]
\[
\therefore\ V=2\pi+\frac{16}3
\]
\end{proof}

\newpage

\section*{第4問}

\begin{proof}[解]
実数$x,y$が$x^2+y^2\le1$をみたすとする。

\textbf{(1)} $s=x+y,\ t=xy$とおく。$x,y$は$z^2-sz+t=0$の2実解だから，判別式条件
\[
s^2-4t\ge0
\]
また$x^2+y^2=s^2-2t\le1$だから
\[
s^2-2t\le1 \iff t\ge\frac{s^2-1}2
\]
2式をまとめて，$(s,t)$の動く範囲は
\[
\frac{s^2-1}2\le t\le\frac{s^2}4
\]
（境界は$t=\frac14s^2$と$t=\frac12(s^2-1)$，交点は$\frac14s^2=\frac12s^2-\frac12\iff s^2=2\iff s=\pm\sqrt2$）。これは$-\sqrt2\le s\le\sqrt2$の範囲の，2つの放物線に挟まれた領域（境界を含む）である。

\textbf{(2)} $f(s,t)=t+ms$とおく。(1)から$\dfrac{s^2-1}2\le t\le\dfrac{s^2}4$（$-\sqrt2\le s\le\sqrt2$）。$s$を固定すると，$f$は$t$の増加関数だから
\[
\frac{s^2-1}2+ms\le f(s,t)\le\frac{s^2}4+ms
\]
左辺を$g(s)$，右辺を$h(s)$とおく（$-\sqrt2\le s\le\sqrt2$）。
\[
g(s)=\frac12(s^2-1)+ms=\frac12(s+m)^2-\frac12m^2-\frac12
\]
\[
h(s)=\frac14s^2+ms=\frac14(s+2m)^2-m^2
\]
$m\ge0$に注意する。

\textbf{最小値について：} $g$は$s=-m$で最小の下に凸な放物線。
\begin{itemize}
\item $-\sqrt2\le-m$（$m\le\sqrt2$）の時：頂点が区間内にあり，$\min g=g(-m)=-\dfrac12m^2-\dfrac12$
\item $-m<-\sqrt2$（$m>\sqrt2$）の時：$g$は区間内で単調増加，$\min g=g(-\sqrt2)=\dfrac12(m-\sqrt2)^2-\dfrac12m^2-\dfrac12=-\sqrt2m+\dfrac12$
\end{itemize}

\textbf{最大値について：} $h$は$s=-2m\le0$で最小をとる下に凸な放物線だから，区間$[-\sqrt2,\sqrt2]$上の最大値は端点でとる。$m\ge0$では頂点$-2m$から遠いのは常に$s=\sqrt2$の側（$|\sqrt2-(-2m)|=\sqrt2+2m\ge|-\sqrt2-(-2m)|=|2m-\sqrt2|$，$m\ge0$で常に成立）だから
\[
\max h=h(\sqrt2)=\frac14(\sqrt2+2m)^2-m^2=\sqrt2m+\frac12
\]

以上をまとめて，もとめる最大値・最小値は
\[
\text{最大値}=\sqrt2\,m+\frac12,\qquad
\text{最小値}=
\begin{cases}
-\dfrac12m^2-\dfrac12 & (0\le m\le\sqrt2)\\[2mm]
-\sqrt2\,m+\dfrac12 & (m\ge\sqrt2)
\end{cases}
\]
\end{proof}

\end{document}
